\documentclass[10pt,a4paper]{amsart}
\usepackage{amsmath,amssymb,amsthm,graphicx}
% Retained minimal preamble for the Stacks Project chapter flat.tex (More on Flatness).
% Verbatim excerpt of the chapter preamble: Xy-pic setup (native preamble lines 14-18) and the native
% theorem declarations with the mathematical macros (native preamble lines 154-217). Copied mechanically;
% nothing added, removed or redefined. The wrapper class and the amsmath/amssymb/amsthm/graphicx base
% packages are supplied by the assembler, so the class/externaldocument part of the native preamble is omitted.
% For commutative diagrams we use Xy-pic
\usepackage[all]{xy}

% We use 2cell for 2-commutative diagrams.
\xyoption{2cell}
% Theorem environments.
%
\theoremstyle{plain}
\newtheorem{theorem}[subsection]{Theorem}
\newtheorem{proposition}[subsection]{Proposition}
\newtheorem{lemma}[subsection]{Lemma}

\theoremstyle{definition}
\newtheorem{definition}[subsection]{Definition}
\newtheorem{example}[subsection]{Example}
\newtheorem{exercise}[subsection]{Exercise}
\newtheorem{situation}[subsection]{Situation}

\theoremstyle{remark}
\newtheorem{remark}[subsection]{Remark}
\newtheorem{remarks}[subsection]{Remarks}

\numberwithin{equation}{subsection}

% Macros
%
\def\lim{\mathop{\mathrm{lim}}\nolimits}
\def\colim{\mathop{\mathrm{colim}}\nolimits}
\def\Spec{\mathop{\mathrm{Spec}}}
\def\Hom{\mathop{\mathrm{Hom}}\nolimits}
\def\Ext{\mathop{\mathrm{Ext}}\nolimits}
\def\SheafHom{\mathop{\mathcal{H}\!\mathit{om}}\nolimits}
\def\SheafExt{\mathop{\mathcal{E}\!\mathit{xt}}\nolimits}
\def\Sch{\mathit{Sch}}
\def\Mor{\mathop{\mathrm{Mor}}\nolimits}
\def\Ob{\mathop{\mathrm{Ob}}\nolimits}
\def\Sh{\mathop{\mathit{Sh}}\nolimits}
\def\NL{\mathop{N\!L}\nolimits}
\def\CH{\mathop{\mathrm{CH}}\nolimits}
\def\proetale{{pro\text{-}\acute{e}tale}}
\def\etale{{\acute{e}tale}}
\def\QCoh{\mathit{QCoh}}
\def\Ker{\mathop{\mathrm{Ker}}}
\def\Im{\mathop{\mathrm{Im}}}
\def\Coker{\mathop{\mathrm{Coker}}}
\def\Coim{\mathop{\mathrm{Coim}}}

% Boxtimes
%
\DeclareMathSymbol{\boxtimes}{\mathbin}{AMSa}{"02}

%
% Macros for moduli stacks/spaces
%
\def\QCohstack{\mathcal{QC}\!\mathit{oh}}
\def\Cohstack{\mathcal{C}\!\mathit{oh}}
\def\Spacesstack{\mathcal{S}\!\mathit{paces}}
\def\Quotfunctor{\mathrm{Quot}}
\def\Hilbfunctor{\mathrm{Hilb}}
\def\Curvesstack{\mathcal{C}\!\mathit{urves}}
\def\Polarizedstack{\mathcal{P}\!\mathit{olarized}}
\def\Complexesstack{\mathcal{C}\!\mathit{omplexes}}
% \Pic is the operator that assigns to X its picard group, usage \Pic(X)
% \Picardstack_{X/B} denotes the Picard stack of X over B
% \Picardfunctor_{X/B} denotes the Picard functor of X over B
\def\Pic{\mathop{\mathrm{Pic}}\nolimits}
\def\Picardstack{\mathcal{P}\!\mathit{ic}}
\def\Picardfunctor{\mathrm{Pic}}
\def\Deformationcategory{\mathcal{D}\!\mathit{ef}}

% --- appended 2026-10-12 (Chief): the project own definitions of the reference/slogan/history
% display environments, copied mechanically from the shared Stacks preamble lines 8-12
% (\usepackage{verbatim} plus the three \newenvironment lines, which the chapter preamble keeps
% verbatim from the project preamble). Nothing added, removed or redefined: these environments
% carry no mathematical content and the project itself typesets them as comments. Without these
% three lines a counted statement whose own text uses them fails with Environment slogan undefined.
\usepackage{verbatim}
\newenvironment{reference}{\comment}{\endcomment}
\newenvironment{slogan}{\comment}{\endcomment}
\newenvironment{history}{\comment}{\endcomment}
\usepackage{cleveref}
\crefname{originallemma}{original lemma}{original lemmas}
\Crefname{originallemma}{Original lemma}{Original lemmas}
\makeatletter
\newlabel{algebra-lemma-cover}{{00EO}{source}{}{sourcepointer-1}{}}
\newlabel{algebra-lemma-cover@cref}{{[originallemma][0][]00EO}{[1][0][]source}{}{}{}}
\makeatother
\crefname{originallemma}{original lemma}{original lemmas}
\Crefname{originallemma}{Original lemma}{Original lemmas}
\makeatletter
\newlabel{algebra-lemma-qc-open}{{00F6}{source}{}{sourcepointer-2}{}}
\newlabel{algebra-lemma-qc-open@cref}{{[originallemma][0][]00F6}{[1][0][]source}{}{}{}}
\makeatother
\crefname{originaltheorem}{original theorem}{original theorems}
\Crefname{originaltheorem}{Original theorem}{Original theorems}
\makeatletter
\newlabel{algebra-theorem-openness-flatness}{{00RC}{source}{}{sourcepointer-3}{}}
\newlabel{algebra-theorem-openness-flatness@cref}{{[originaltheorem][0][]00RC}{[1][0][]source}{}{}{}}
\makeatother
\crefname{originallemma}{original lemma}{original lemmas}
\Crefname{originallemma}{Original lemma}{Original lemmas}
\makeatletter
\newlabel{lemma-finite-type-flat-at-point-local}{{05I7}{source}{}{sourcepointer-4}{}}
\newlabel{lemma-finite-type-flat-at-point-local@cref}{{[originallemma][0][]05I7}{[1][0][]source}{}{}{}}
\makeatother
\title{A finite type module flat over the base and finitely presented over a scheme theoretically dense open is finitely presented}
\author{The Stacks Project Authors}
\date{stacks-project:flat.tex:local-snapshot-20261011; GNU FDL 1.2 or later}
\begin{document}
\maketitle
\noindent\textit{Editorial note.} All delivered mathematical text is the author's own, copied from the pinned Stacks Project chapter file flat.tex; no statement or proof step was rewritten, reordered or removed.\par\medskip
\noindent\textit{Editorial note.} Preamble provenance: the chapter's own retained minimal preamble (sources/stacks-flat/retained-preamble.tex, lines 1-85) is a verbatim excerpt of the native Xy-pic setup and of the native theorem declarations and mathematical macros, followed by the project's own reference/slogan/history display-environment definitions copied mechanically from the shared preamble; the wrapper class and the amsmath/amssymb/amsthm/graphicx base packages are supplied by the assembler, and extra\_packages is left empty because everything needed is already in the retained preamble. No macro was invented or redefined.\par\medskip
\noindent\textit{Editorial note.} Same-chapter pointers are located at their \textbackslash{}label lines inside the pinned chapter file; cross-chapter pointers are located in the project's own tag list staged in the snapshot (sources/stacks-chow/tags.txt), at the verbatim lines carrying their canonical tags.\par\medskip
\noindent\textit{Editorial note.} The counted statement contains the author's own slogan display environment at source lines 2944-2947, which the project typesets as a comment; the retained chapter preamble at sources/stacks-flat/retained-preamble.tex supplies its definition (lines 1-85 of that file), so the statement compiles.\par\medskip
\noindent\textit{Allowed standard prerequisites (editorial).}\begin{itemize}
\item A quasi-compact open of an affine scheme is a finite union of principal opens. Algebra, label algebra-lemma-qc-open, invoked at proof line 2965. Writes the dense open U as a finite union of distinguished opens D(f\_i), which is what makes the injectivity of A -> product of the localisations available.
\item Finiteness of presentation at a point for a flat finite type module, and openness of flatness. Same chapter flat.tex, label lemma-finite-type-flat-at-point-local (proof line 2970); Algebra, label algebra-theorem-openness-flatness (proof line 2981). Gives that the localisation of the module at the chosen prime is finitely presented over the local ring and, after constructing a model that agrees with it there, that the model is flat over the base on a neighbourhood of the point.
\item A local presentation argument, and the local nature of finite presentation. Algebra, labels algebra-lemma-cover (proof line 2996) and algebra-lemma-extension (used in the same section). The cover lemma lets a local finite presentation statement be concluded globally, which is how the proof finishes from the affine neighbourhood found at the chosen point.
\end{itemize}\medskip
\ifdefined\hypertarget\hypertarget{sourcepointer-1}{}\fi
\noindent\textit{Source reference.} Original lemma 00EO: sources/stacks-chow/tags.txt, source lines 542--542. Navigation pointer for \textbackslash{}ref\{algebra-lemma-cover\} used inside the counted span (line 2996); the argument that consumes it is retained in the span.\par\medskip
\ifdefined\hypertarget\hypertarget{sourcepointer-2}{}\fi
\noindent\textit{Source reference.} Original lemma 00F6: sources/stacks-chow/tags.txt, source lines 560--560. Navigation pointer for \textbackslash{}ref\{algebra-lemma-qc-open\} used inside the counted span (line 2965); the argument that consumes it is retained in the span.\par\medskip
\ifdefined\hypertarget\hypertarget{sourcepointer-3}{}\fi
\noindent\textit{Source reference.} Original theorem 00RC: sources/stacks-chow/tags.txt, source lines 998--998. Navigation pointer for \textbackslash{}ref\{algebra-theorem-openness-flatness\} used inside the counted span (line 2981); the argument that consumes it is retained in the span.\par\medskip
\ifdefined\hypertarget\hypertarget{sourcepointer-4}{}\fi
\noindent\textit{Source reference.} Original lemma 05I7: sources/stacks-flat/pinned-original-source.tex, source lines 2879--2879. Navigation pointer for \textbackslash{}ref\{lemma-finite-type-flat-at-point-local\} used inside the counted span (line 2970); the argument that consumes it is retained in the span.\par\medskip
\begin{lemma}
\label{lemma-flat-finite-type-finitely-presented-over-dense-open}
\begin{slogan}
$S$-flat and finite type extensions of finitely presented modules
on a (good) open are also $X$-finitely presented.
\end{slogan}
Let $f : X \to S$ be a morphism of schemes. Let $\mathcal{F}$ be a
quasi-coherent $\mathcal{O}_X$-module. Let $U \subset S$ be open.
Assume
\begin{enumerate}
\item $f$ is locally of finite presentation,
\item $\mathcal{F}$ is of finite type and flat over $S$,
\item $U \subset S$ is retrocompact and scheme theoretically dense,
\item $\mathcal{F}|_{f^{-1}U}$ is of finite presentation.
\end{enumerate}
Then $\mathcal{F}$ is of finite presentation.
\end{lemma}

\begin{proof}
The problem is local on $X$ and $S$, hence we may assume $X$ and $S$ affine.
Write $S = \Spec(A)$ and $X = \Spec(B)$. Let $N$ be a finite $B$-module such
that $\mathcal{F}$ is the quasi-coherent sheaf associated to $N$.
We have $U = D(f_1) \cup \ldots \cup D(f_n)$ for some $f_i \in A$, see
Algebra, Lemma \ref{algebra-lemma-qc-open}.
As $U$ is schematically dense the map
$A \to A_{f_1} \times \ldots \times A_{f_n}$ is injective.
Pick a prime $\mathfrak q \subset B$ lying over $\mathfrak p \subset A$
corresponding to $x \in X$ mapping to $s \in S$.
By Lemma \ref{lemma-finite-type-flat-at-point-local}
the module $N_\mathfrak q$ is of finite presentation over $B_\mathfrak q$.
Choose a surjection $\varphi : B^{\oplus m} \to N$ of $B$-modules.
Choose $k_1, \ldots, k_t \in \Ker(\varphi)$ and
set $N' = B^{\oplus m}/\sum Bk_j$. There is a canonical surjection
$N' \to N$ and $N$ is the filtered colimit of the $B$-modules $N'$
constructed in this manner. Thus we see that we can choose
$k_1, \ldots, k_t$ such that (a) $N'_{f_i} \cong N_{f_i}$, $i = 1, \ldots, n$
and (b) $N'_\mathfrak q \cong N_\mathfrak q$.
This in particular implies that $N'_\mathfrak q$ is flat over $A$.
By openness of flatness, see
Algebra, Theorem \ref{algebra-theorem-openness-flatness}
we conclude that there exists a $g \in B$, $g \not \in \mathfrak q$
such that $N'_g$ is flat over $A$. Consider the commutative diagram
$$
\xymatrix{
N'_g \ar[r] \ar[d] & N_g \ar[d] \\
\prod N'_{gf_i} \ar[r] & \prod N_{gf_i}
}
$$
The bottom arrow is an isomorphism by choice of $k_1, \ldots, k_t$.
The left vertical arrow is an injective map as
$A \to \prod A_{f_i}$ is injective and $N'_g$ is flat over $A$.
Hence the top horizontal arrow is injective, hence an isomorphism.
This proves that $N_g$ is of finite presentation over $B_g$.
We conclude by applying
Algebra, Lemma \ref{algebra-lemma-cover}.
\end{proof}

\end{document}
